\chapter{Gestion de Version avec Git et GitHub}

Dans le cadre de ce projet, Git a été utilisé comme outil principal de gestion de version. Cette section met en évidence les différentes étapes de suivi et de collaboration, notamment les commits clairs, les pull requests, et les tags.

\section{Github Repository}
Le projet est hébergé sur une plateforme GitHub, qui permet une gestion efficace du code source et une collaboration fluide entre les membres de l’équipe. Le dépôt GitHub contient tout le code du projet, organisé en plusieurs dossiers et fichiers pour une meilleure structure.

\section{Network}
L’onglet \texttt{Network} de GitHub fournit une visualisation graphique de l’historique des branches et des fusions. Cet outil est essentiel pour suivre l’évolution du projet et comprendre les interactions entre les différentes branches.

\begin{figure}[H]
    \centering
    %\includegraphics[width=1\textwidth]{fig/network.png}
    \caption{Vue du réseau des branches et des commits.}
    \label{fig:network}
\end{figure}
